\chapter{Quantum-tunnelling relaxation prediction (menu 46)}
\label{ch:menu46}
\menulabel{46}

\section{Purpose}

Predicts the tunnelling relaxation time $\tau_{\mathrm{QTM}}$ and the
effective demagnetization barrier $U_{\mathrm{eff}}(T)$ of a Kramers
single-ion magnet from its ab initio parameters, through two published
models over the SINGLE\_ANISO data chain of menu 36.

\section{What you need}

An ORCA output with a SINGLE\_ANISO section; the $B_{\mathrm{ave}}$ field
scale (Enter $=$ 20.0~mT); optionally a neighbour table for the
spin-dipolar model (\code{dx dy dz mx my mz} per line -- position in
Angstrom and moment in Bohr magnetons, both in the central ion's
principal-g frame), plus optional dilution concentrations.

\section{Walkthrough}

The example reads the CO$^{+}$ fixture through both routes: the
equivalent-Zeeman table (near-isotropic $g$ gives fast QTM,
$\tau = 1.547\times10^{-9}$~s; the excited doublet at
29361~cm$^{-1}$ lies far above the grid, and the report says so), then
with the constructed twelve-centre neighbour probe through the dipolar
model and the seeded dilution table.

\lstinputlisting[style=fbkcode, firstline=59, lastline=69,
  title={The menu-46 example (the doublet table and the ground-state
  tunnelling time)}]
  {examples/expected/m46_qtm.txt}

\lstinputlisting[style=fbkcode, firstline=179, lastline=190,
  title={The spin-dipolar block with the dilution medians (the constructed
  neighbour probe)}]
  {examples/expected/m46_qtm.txt}

\section{What you get}

The report \code{<output>.qtm.fbk.md}: the per-doublet table (energy,
$g_{XY}$, $g_Z$, $\tau$), the ground-state $\tau_{\mathrm{QTM}}$, the
$U_{\mathrm{eff}}(T)$ table with per-doublet contributions, and -- when a
neighbour table is given -- the dipolar block ($\sigma_r/\sigma_i$,
$\langle|E_{sf}|\rangle$, $\tau$) and the dilution medians.

\section{Conventions, cross-checks and boundaries (measured)}

\begin{readbox}
\begin{itemize}
  \item \textbf{The equivalent-Zeeman regression is exact to the print
    precision}: the shipped 18-complex literature table ($g$ from the
    JPCL 2018 SI, $\log\tau$ from the PCCP 2020 Table~1, two independent
    data chains of the same model) reproduces to a maximum deviation of
    0.006 in $\log_{10}(\tau)$ at $B_{\mathrm{ave}} = 20$~mT with CODATA
    constants.
  \item \textbf{The dipolar Fermi-golden-rule constant carries the
    source's one-unit density factor $1/k_B$}: the published
    $\langle|E_{sf}|\rangle$ values of three JPCL-SI complexes (kelvin)
    reproduce the published tunnelling times through
    $k = (2\pi k_B/\hbar)\langle|E_{sf}|\rangle_K^2$ -- NAFMIT
    $8.671\times10^{-9}$~K gives $\log_{10}\tau = 3.906$ against the
    published 3.91, BAJSIQ $-4.586$ against $-4.59$.
  \item \textbf{Axial convention}: the largest principal $g$ is taken as
    the axial component (the formula pair enters $g_{XY}^2$
    symmetrically); the models are zero-field, Kramers-ion and
    single-centre.
  \item \textbf{$B_{\mathrm{ave}}$ is an empirical field scale} (20~mT in
    the source; the paper states it can be treated as an empirical
    parameter) -- the value is a menu question.
  \item \textbf{These predictions are absolute-value estimates}; menu
    36's reading of an experimental-style plateau is a separate,
    data-side quantity -- keep the two apart when quoting them side by
    side.
  \item \textbf{Scope}: the neighbour geometry enters as a
    table (crystal-structure parsing sits outside this menu).  A recorded open
    point: the literature reports a roughly quadratic
    $\log\tau$-$\log x$ dilution slope ($-1.93$), while the independent
    Bernoulli construction here approaches the analytic expectation
    ($-1$) for large neighbour counts; the difference is recorded, not
    tuned away.
\end{itemize}
\end{readbox}
